\section{Abstract}

The report explored the application of convolutional neural networks (CNNs) for the classification of steady-state visually evoked popetentials (SSVEPs) in brain-computer interfaces (BCIs).The main goal was to enhance classification accuracy and demonstrate the potential of deep learning methodologies in processing neural signal data.

Using a benchmark dataset consisting of 64-channel EEG recordings from 35 participants engaged in a cue-guided target selection task. Firstly, we employed signal processing techniques and feature extraction methods, in order to perform profiling and visualize the data. These steps included band-pass filtering, power spectral density (PSD) analysis, and wavelet transforms.

Finally, an end-to-end CNN-based classifier was designed and implemented in PyTorch. The architecture incorporated temporal and spatial convolutions, dropout layers, batch normalization, and data augmentation strategies to improve performance and generalization. Hyperparameter optimization and ablation studies were conducted to validate the effectiveness of various model components.

With the deep learning approach, we obtained validation accuracies of 97.50\% and 96.67\% for the 2 proposed architectures, with satisfactory training curves, demonstrating very little overfitting, good discrimination between classes and good generalization to unseen data. These results are very satisfactory, and demonstrate the incredible ability of CNNs to perform end-to-end classification on EEG signals, without prior domain knowledge or feature extraction.
